% TOP_PROBLEM: {"id": 1987, "title": "Erdős Problem #196", "category_id": 2, "category": {"id": 2, "name": "combinatorics", "display_name": "Combinatorics", "description": "Counting problems, graph theory, discrete structures.", "slug": "combinatorics", "order_index": 2, "created_at": "2026-07-31T15:26:25.671Z"}, "status": "open"}
% FIRST_POSTED: null
% ENTRY_KIND: statement_only
% Imported from: unsolvedmath_additional_500.tex; UnsolvedMath ID 1987
% !TeX program = lualatex
% Compile twice: lualatex 1987.tex
% Self-contained: no external bibliography, images, or \input files.
% Numeric section labels and visible numbers are original UnsolvedMath IDs.
\documentclass[11pt,a4paper]{article}
\usepackage[margin=24mm,headheight=14pt]{geometry}
\usepackage{fontspec}
\setmainfont{Latin Modern Roman}
\setsansfont{Latin Modern Sans}
\setmonofont{Latin Modern Mono}
\usepackage{amsmath,amssymb,mathtools}
\usepackage{unicode-math}
\setmathfont{Latin Modern Math}
\usepackage{microtype}
\usepackage{enumitem,xurl,needspace}
\usepackage[dvipsnames]{xcolor}
\usepackage{hyperref}
\hypersetup{unicode=true,colorlinks=true,linkcolor=MidnightBlue,urlcolor=MidnightBlue,
 pdftitle={UnsolvedMath Problem 1987},pdfauthor={Source-annotated compilation},bookmarksnumbered=true}
\usepackage{bookmark,tocloft,titlesec,fancyhdr}
\setlength{\cftsecnumwidth}{4.2em}
\cftsetpnumwidth{2.5em}
\cftsetrmarg{4em}
\titleformat{\section}{\Large\bfseries\raggedright}{\thesection}{0.7em}{}
\pagestyle{fancy}\fancyhf{}
\fancyhead[L]{\small\sffamily Additional UnsolvedMath statements}
\fancyhead[R]{\small Original catalogue IDs}
\fancyfoot[C]{\thepage}
\setlength{\parindent}{0pt}
\setlength{\parskip}{5pt plus 1pt}
\setlength{\emergencystretch}{3em}
\setcounter{tocdepth}{1}\setcounter{secnumdepth}{1}
\setlist[itemize]{leftmargin=3em,itemsep=4pt,topsep=4pt,parsep=0pt}
\allowdisplaybreaks
\newcommand{\N}{\mathbb{N}}
\newcommand{\Z}{\mathbb{Z}}
\newcommand{\Q}{\mathbb{Q}}
\newcommand{\R}{\mathbb{R}}
\newcommand{\C}{\mathbb{C}}
\newcommand{\F}{\mathbb{F}}
\newcommand{\A}{\mathbb{A}}
\newcommand{\PP}{\mathbb{P}}
\newcommand{\E}{\mathbb{E}}
\newcommand{\eps}{\varepsilon}
\newcommand{\dd}{\,\mathrm d}
\newcommand{\sref}[2]{\hyperlink{source:#1:#2}{[S#2]}}
\newif\ifshowcatalogue
\showcataloguetrue
\newcommand{\cataloguescope}[1]{\ifshowcatalogue
 \paragraph{Statement recorded in the supplied source.}
 #1\fi}
\begin{document}
% UnsolvedMath ID 1987; source catalogue code EP-196

\setcounter{section}{1986}

\section{Four-Term Arithmetic Progressions in Permutations}\label{1987}

{\small\sffamily UnsolvedMath ID 1987 · Combinatorics}

\subsection{Definitions and mathematical statement}

\paragraph{Shared notation.}Unless a section says otherwise, $\N=\{1,2,\ldots\}$,
$\Z$ is the ring of integers, $\Q,\R,\C$ are the rational, real, and complex
fields, $[N]=\{1,\ldots,\lfloor N\rfloor\}$, and $\log$ is the natural
logarithm. For real functions of a parameter tending to infinity,
$f\ll g$ and $f=O(g)$ mean $|f|\leq Cg$ eventually for some fixed $C$;
$f\gg g$ reverses this comparison; $f=o(g)$ means $f/g\to0$;
$f\sim g$ means $f/g\to1$; and $f\asymp g$ means both order bounds.
A subscript on $O$ or $\ll$ specifies allowed constant dependence.
The expression $n^{o(1)}$ in an upper bound means growth smaller than
$n^\epsilon$ for each fixed $\epsilon>0$. Local definitions govern any
exception, finite-field convention, or use of the same symbol for another
object. An asymptotic claim is not automatically an effective algorithm.



A permutation of $\N$ is a bijection $x:\N\to\N$. The entries need not be at consecutive positions. They form a nonconstant arithmetic progression when read in the specified increasing or decreasing order of their indices. \sref{1987}{1} \sref{1987}{2}

\paragraph{Statement recorded in the supplied source.}
Must every permutation of $\mathbb{N}$ contain a monotone 4-term arithmetic progression? In other words, given a permutation $x$ of $\mathbb{N}$ must there be indices with either $i<j<k<l$ or $i>j>k>l$ such that $x_i,x_j,x_k,x_l$ are an arithmetic progression? \sref{1987}{1}

\subsection{Short English statement}

Does every ordering of the positive integers contain four entries, appearing in forward or reverse order, that form an arithmetic progression?

\subsection{Sources}

{\small

\begin{itemize}

\item[\hypertarget{source:1987:1}{S1}] ULAM.AI, \emph{UnsolvedMath}, supplied \texttt{problems.json}, record 1987 (EP-196), statement and background. The embedded literature-review date is 2026-08-17; that is the archive’s review date, not a fresh certification. Dataset landing page: \url{https://huggingface.co/datasets/ulamai/UnsolvedMath}.

\item[\hypertarget{source:1987:2}{S2}] T. F. Bloom, \emph{Erdős Problems}, Problem 196, \url{https://www.erdosproblems.com/196}. Reference imported from the supplied record; the live problem page was not independently checked for this compilation.

\item[\hypertarget{source:1987:3}{S3}] J. A. Davis, R. C. Entringer, R. L. Graham, and G. J. Simmons, On permutations containing no long arithmetic progressions, Acta Arithmetica 34 (1977/78), 81-90. \url{https://eudml.org/doc/205322}. Bibliographic information imported from the supplied record.

\item[\hypertarget{source:1987:4}{S4}] Timothy D. LeSaulnier and Sujith Vijay, On permutations avoiding arithmetic progressions, Discrete Mathematics 311 (2011), 205-207. \url{https://doi.org/10.1016/j.disc.2010.10.006}. Bibliographic information imported from the supplied record.

\item[\hypertarget{source:1987:5}{S5}] Sarosh Adenwalla, A Generalisation of a Result on Monotone Arithmetic Progressions in Permutations of the Positive Integers, arXiv:2302.09662. \url{https://arxiv.org/abs/2302.09662}. Bibliographic information imported from the supplied record.

\end{itemize}

}

\label{end:1987}
\end{document}
